The phrase ``$2^n-1,3^n-1,\dots,h(n)^n-1$ are mutually coprime'' is ambiguous: it can mean
\begin{itemize}[leftmargin=2em]
\item \emph{pairwise} coprime (every pair has gcd $1$), or
\item \emph{collectively} coprime (their \emph{common} gcd is $1$).
\end{itemize}
If interpreted pairwise, the definition would be degenerate (already a single number is vacuously pairwise coprime), and the stated fact ``$h(n)=n+1$ iff $n+1$ is prime'' would fail.
The only interpretation consistent with the rest of the prompt is the \emph{collective} one.

\begin{definition}[$h(n)$]
For $n\ge 1$, let $h(n)$ be the smallest integer $m\ge 2$ such that
\[
\gcd\bigl(2^n-1,3^n-1,\dots,m^n-1\bigr)=1.
\]
(Equivalently: the integers $2^n-1,\dots,m^n-1$ have no common prime divisor.)
\end{definition}

\paragraph{Questions (Erd\H{o}s).}
\begin{enumerate}[label=(\alph*),leftmargin=2.2em]
\item For each prime $p$, does the natural density $\delta_p$ of integers $n$ with $h(n)=p$ exist?
\item Is $\liminf_{n\to\infty} h(n)=\infty$?
\item If $p$ is the greatest prime such that $p-1\mid n$ and $p>n^{\varepsilon}$, must $h(n)=p$?
\end{enumerate}
